%% kapa: exam=1 round=35 year=2024 period=1 form=A subject=real_estate subject_ko=부동산학원론 seq=35 number=115
%% source: Q-Net 감정평가사 기출문제 (cst00309) artlSeq=5213229 file=raw/1008/5213229/1교시 문제지 원본.pdf
%% conversion: kordoc 4.24.51 -> tools/gen1.py (auto); review=proofread-sonnet; flags=none
\begin{question}{115}
부동산 마케팅활동에 관한 설명으로 옳지 않은 것은?
\begin{choices}
  \item 시장세분화란 부동산시장에서 마케팅활동을 수행하기 위하여 구매자의 집단을 세분화하는 것이다.
  \item 세분시장은 그 규모와 구매력 등의 특성이 측정될 수 있어야 한다.
  \item 세분시장은 개념적으로 구분될 수 있으며 마케팅 믹스 요소에 대해 동일하게 반응한다.
  \item 표적시장이란 세분화된 시장 중 가장 효과적인 성과가 기대되어 마케팅활동의 수행대상이 되는 시장을 말한다.
  \item 포지셔닝은 표적시장에서 고객의 욕구를 파악하여 경쟁제품과 차별화된 자사제품의 개념을 정해 이를 소비자의 지각 속에 적절히 위치시키는 것이다.
\end{choices}
\end{question}
\answer{3}
